%========================================================================================
\section{Sécurité}
\label{sec:securite}
%========================================================================================

\subsection{Authentification par JWT}

L'API est sans état : chaque requête porte un jeton JWT~\cite{rfc7519}, émis par SimpleJWT~\cite{simplejwt}. Le jeton d'accès expire après 15~minutes, ce qui limite l'effet d'un vol. Le jeton de rafraîchissement dure un jour et ne sert qu'une fois : chaque rafraîchissement en émet un nouveau et place l'ancien sur une liste noire. La déconnexion révoque le jeton fourni, à condition qu'il appartienne à l'utilisateur connecté.

\begin{lstlisting}[language=Python,caption={Paramètres JWT (\texttt{config/settings/base.py}).},label={lst:jwt}]
SIMPLE_JWT = {
    "ACCESS_TOKEN_LIFETIME": timedelta(minutes=15),
    "REFRESH_TOKEN_LIFETIME": timedelta(days=1),
    # Chaque rafraîchissement émet un nouveau refresh et met l'ancien sur liste noire.
    "ROTATE_REFRESH_TOKENS": True,
    "BLACKLIST_AFTER_ROTATION": True,
    "UPDATE_LAST_LOGIN": True,
    "AUTH_HEADER_TYPES": ("Bearer",),
}
\end{lstlisting}

Les mots de passe passent les quatre validateurs de Django : longueur minimale, similarité avec le profil, mots de passe courants, valeurs entièrement numériques.

\subsection{Permissions par rôle}

DRF refuse par défaut toute requête anonyme (\code{IsAuthenticated}). Trois permissions objet (\code{apps/core/permissions.py}) traduisent ensuite les rôles : \code{IsTripMember}, \code{IsTripEditor} et \code{IsTripOwner}. Elles acceptent un voyage ou tout objet qui s'y rattache ; une activité remonte à son voyage par \code{activity.destination.trip}.

\begin{table}[H]
  \centering
  \small
  \caption{Matrice d'accès à un voyage (codes HTTP).}
  \label{tab:matrice}
  \begin{tabularx}{\linewidth}{@{}X c c c c@{}}
    \toprule
    \textbf{Opération} & \textbf{Propriétaire} & \textbf{Éditeur} & \textbf{Lecteur} & \textbf{Non participant} \\
    \midrule
    Lire le voyage, \code{stats}, \code{clone} & 200 / 201 & 200 / 201 & 200 / 201 & 404 \\
    Modifier le voyage, ses destinations, ses activités & 200 & 200 & 403 & 404 \\
    Supprimer le voyage & 204 & 403 & 403 & 404 \\
    Gérer les participants & 201 / 200 / 204 & 403 & 403 & 404 \\
    \bottomrule
  \end{tabularx}
\end{table}

\subsection{Contrôle d'accès au niveau objet}

Le risque principal d'une API est le contrôle d'accès défaillant au niveau objet (\emph{Broken Object Level Authorization}, API1 de l'OWASP~\cite{owasp}) : modifier l'identifiant dans l'URL pour atteindre l'objet d'autrui. TravelHub s'en protège par construction. Chaque vue, chaque resolver et chaque filtre part de \code{visible_to(user)}. Un voyage d'autrui est absent du queryset, et la réponse est 404. Le même principe s'applique au parent d'une création : une destination créée sur le voyage d'autrui donne 404, sur un voyage visible sans droit d'écriture 403.

\subsection{Configuration de production}

Le module \code{config.settings.prod} active les protections du transport et du navigateur : redirection HTTPS, HSTS de 30 jours, cookies \code{Secure}, \code{X-Frame-Options: DENY}, \code{nosniff}, politique de référent et COOP. L'API navigable de DRF y est désactivée : seul le JSON est servi. GraphiQL et l'introspection y sont fermés. La commande suivante, exécutée par un test, ne produit aucun avertissement~\cite{djangodeploy} :

\begin{lstlisting}[language={},caption={Vérification de la configuration de production.},label={lst:deploy}]
python manage.py check --deploy --settings=config.settings.prod
\end{lstlisting}

\paragraph{CORS.} Seules les origines listées dans \code{CORS_ALLOWED_ORIGINS} sont acceptées, sans transmission des identifiants. La liste est vide par défaut ; le TP2 y ajoutera l'origine du front-end.

\paragraph{Secrets.} La clé secrète et l'accès à la base sont lus dans l'environnement. Le fichier \code{.env} est ignoré par Git ; \code{.env.example} ne contient que des valeurs fictives. Un test vérifie que \code{SECRET_KEY} provient bien de l'environnement. Les journaux ne contiennent ni corps de requête, ni en-têtes, ni chaîne de requête. En filet de sécurité, le formateur masque tout champ dont le nom évoque un secret, ainsi que tout JWT présent dans un message.

\subsection{OWASP API Security Top 10}

Le tableau~\ref{tab:owasp} résume la couverture des dix risques de l'édition 2023.

\begin{table}[H]
  \centering
  \small
  \caption{Mesures face à l'OWASP API Security Top 10 (2023).}
  \label{tab:owasp}
  \begin{tabularx}{\linewidth}{@{}>{\raggedright\arraybackslash}p{4.2cm} X@{}}
    \toprule
    \textbf{Risque} & \textbf{Mesures} \\
    \midrule
    API1 Object Level Authorization & \code{visible_to()} partout ; 404 plutôt que 403 ; parent cherché parmi les objets visibles \\
    API2 Authentication & Accès de 15 min, rotation et liste noire, validateurs de mot de passe, cookie de session ignoré par GraphQL \\
    API3 Object Property Level Authorization & Listes blanches de champs ; \code{owner}, totaux et horodatages en lecture seule ; pas de \code{is_staff} à l'inscription \\
    API4 Resource Consumption & 100 requêtes/h en anonyme, 1000/h par utilisateur ; pagination ; profondeur GraphQL limitée ; délai de 5 s vers l'extérieur \\
    API5 Function Level Authorization & Permissions par rôle ; seul le propriétaire gère les participants \\
    API6 Sensitive Business Flows & Inscription, connexion et clonage soumis aux limites de requêtes \\
    API7 Server Side Request Forgery & URL externes constantes ; ville en paramètre, code pays en variable GraphQL \\
    API8 Security Misconfiguration & \code{check --deploy} propre ; en-têtes de sécurité ; erreurs JSON sans trace \\
    API9 Inventory Management & Une seule version, décrite par OpenAPI et par la collection Postman \\
    API10 Unsafe Consumption of APIs & Réponses externes validées et converties en \code{dataclasses} ; erreurs 502, 503, 504 \\
    \bottomrule
  \end{tabularx}
\end{table}

Chaque ligne est couverte par des tests dédiés. La matrice d'accès croisés de \code{test_object_level_authorization.py} vérifie, pour chaque point d'accès REST et chaque opération GraphQL, qu'un non-participant, un lecteur et un anonyme sont refusés, et que la base reste inchangée après le refus.

\subsection{Risques résiduels}

Trois risques sont connus et acceptés.
\begin{itemize}
  \item L'inscription et l'ajout d'un participant indiquent si un courriel est inscrit. L'énumération reste possible, freinée par les limites de requêtes.
  \item GraphQL limite la profondeur des requêtes, mais pas leur largeur (alias multiples).
  \item La documentation OpenAPI reste publique en production, alors que l'introspection GraphQL y est fermée.
\end{itemize}
